\expandafter\ifx\csname CAPETINCLUDEMODE\endcsname\relax
\documentclass[11pt,a4paper]{article}
% preamble.tex — packages et commandes communes a toutes les fiches CAPET SII IE
\input{glyphtounicode}
\pdfgentounicode=1
\usepackage{cmap}
\usepackage[T1]{fontenc}
\usepackage[utf8]{inputenc}
\usepackage[french]{babel}
\usepackage[a4paper, top=1.6cm, bottom=1.6cm, left=1.4cm, right=1.4cm]{geometry}
\setlength{\headheight}{14.5pt}
\usepackage{amsmath, amssymb, amsfonts}
\usepackage{siunitx}
\sisetup{locale=FR, output-decimal-marker={,}, inter-unit-product=\ensuremath{{}\cdot{}}}
\usepackage{array, booktabs, multirow, tabularx}
\usepackage{xcolor}
\usepackage[most, breakable]{tcolorbox}
\usepackage{multicol}
\usepackage{enumitem}
\usepackage{mathtools}
\usepackage{lmodern}
\usepackage{microtype}
\usepackage{fancyhdr}
\usepackage{tikz}
\usetikzlibrary{positioning,arrows.meta}
\usepackage{lastpage}

% ── Couleurs par domaine ──────────────────────────────────────────────
\definecolor{colElec}{RGB}{26, 82, 148}       % Électrotechnique — bleu foncé
\definecolor{colElecLight}{RGB}{210, 225, 245} % fond bleu clair
\definecolor{colAuto}{RGB}{34, 120, 66}        % Automatique — vert
\definecolor{colMeca}{RGB}{185, 90, 20}        % Mécanique — orange
\definecolor{colTherm}{RGB}{160, 40, 40}       % Thermique — rouge
\definecolor{colQT}{RGB}{232, 244, 233}        % fond question type (vert très léger)
\definecolor{colJury}{RGB}{255, 243, 205}      % fond jury (jaune ambre)
\definecolor{colJuryFrame}{RGB}{200, 130, 0}   % bordure jury
\definecolor{colGrayLight}{RGB}{245, 245, 247} % fond gris neutre
\definecolor{colSeparator}{RGB}{180, 180, 180} % gris séparateur

% ── Style tcolorbox de section principale ────────────────────────────
\tcbset{
  ficheSection/.style={
    breakable,
    enhanced,
    colback=#1!7!white,
    colframe=#1,
    colbacktitle=#1,
    coltitle=white,
    fonttitle=\bfseries\normalsize,
    left=4pt, right=4pt, top=3pt, bottom=4pt,
    before skip=6pt, after skip=6pt,
  }
}

% ── Style encadré question type ──────────────────────────────────────
\tcbset{
  questiontype/.style={
    breakable,
    enhanced,
    colback=colQT,
    colframe=colElec!55!black,
    colbacktitle=colElec!55!black,
    coltitle=white,
    fonttitle=\bfseries\small,
    left=4pt, right=4pt, top=3pt, bottom=3pt,
    before skip=5pt, after skip=5pt,
  }
}

% ── Style encadré jury ───────────────────────────────────────────────
\tcbset{
  jury/.style={
    breakable,
    enhanced,
    colback=colJury,
    colframe=colJuryFrame,
    fonttitle=\bfseries\normalsize,
    left=4pt, right=4pt, top=3pt, bottom=3pt,
  }
}

% ── Mise en page ─────────────────────────────────────────────────────
\pagestyle{fancy}
\fancyhf{}
\renewcommand{\headrulewidth}{0.4pt}
\renewcommand{\footrulewidth}{0.3pt}

\setlist[itemize]{leftmargin=1.2em, itemsep=1pt, topsep=2pt, parsep=0pt}
\setlist[enumerate]{leftmargin=1.2em, itemsep=3pt, topsep=2pt, parsep=0pt}

\setlength{\parskip}{2pt}
\setlength{\columnsep}{10pt}

\def\CAPETBEGINFILE{\begin{document}}
\def\CAPETENDFILE{\end{document}}
\else
\def\CAPETBEGINFILE{}
\def\CAPETENDFILE{}
\fi

\fancyhead[L]{\textbf{\textcolor{colElec}{T18}} - Circuits electroniques - composants (diode, transistor)}
\fancyhead[R]{\small\textit{CAPET SII IE - Session 2026}}
\fancyfoot[C]{\small\thepage\,/\,\pageref{LastPage}}

\CAPETBEGINFILE

% -----------------------------------------------------------------------
%  BLOC EN-TeTE
% -----------------------------------------------------------------------

\begin{tcolorbox}[
  enhanced, breakable,
  colback=colElec,
  colframe=colElec,
  left=6pt, right=6pt, top=5pt, bottom=5pt,
  before skip=0pt, after skip=4pt
]
  {\color{white}\bfseries\large T18 \quad Circuits electroniques - composants (diode, transistor)}\\[2pt]
  {\color{colElec!30!white}\small
    \textbf{Domaine :} electronique - Composants et convertisseurs \quad
    \textbf{Presence :} \textbf{7/13} sujets (2013-2026) - theme \emph{frequent}
  }\\[1pt]
  {\color{colElec!20!white}\footnotesize
    Sous-themes : diode (modele ideal, seuil, droite de charge) - transistor BJT (3 regimes) - MOSFET et IGBT - hacheur abaisseur (courant inductance, ondulation) - diode de roue libre - rapport cyclique
  }
\end{tcolorbox}

% -----------------------------------------------------------------------
%  1. RAPPELS DE COURS
% -----------------------------------------------------------------------

\begin{tcolorbox}[
  ficheSection=colElec,
  title={1. Rappels de cours},
]

% -- 1.1 Diode --------------------------------------------------------
\textbf{1.1 La diode}

\smallskip
\begin{center}
\renewcommand{\arraystretch}{1.4}
\begin{tabular}{lll}
\toprule
\textbf{Modele} & \textbf{etat conducteur ($V_{AK}>0$)} & \textbf{etat bloque ($V_{AK}<0$)} \\
\midrule
Ideal & $V_{AK} = 0$, $i \geq 0$ & $i = 0$ \\
Avec seuil & $V_{AK} = V_{seuil}$ (Si : 0{,}6\,V) & $i = 0$ \\
\bottomrule
\end{tabular}
\end{center}
\textbf{Droite de charge} pour un circuit $E = R\,i + V_{AK}$ : pente $= -1/R$, intersection axe $i$ en $E/R$ et axe $V$ en $E$. Le \emph{point de fonctionnement} Q est l'intersection avec la caracteristique $i(V_{AK})$.

\textbf{Diode de roue libre} : placee en anti-parallele avec une charge inductive. - l'ouverture de l'interrupteur, elle assure la continuite du courant dans l'inductance et protege le composant actif contre les surtensions $L\,di/dt$.

\medskip\hrule\vspace{4pt}

% -- 1.2 Transistor bipolaire (BJT) -----------------------------------
\textbf{1.2 Transistor bipolaire (BJT) - trois regimes}

\smallskip
\begin{center}
\renewcommand{\arraystretch}{1.3}
\begin{tabular}{llll}
\toprule
\textbf{Regime} & \textbf{Condition} & \textbf{Courant $I_C$} & \textbf{Usage} \\
\midrule
Bloque & $V_{BE} < 0{,}6\,\mathrm{V}$ & $\approx 0$ & Interrupteur ouvert \\
Actif (lineaire) & $V_{BE} \approx 0{,}6\,\mathrm{V}$, $V_{CE} > V_{CE(sat)}$ & $\beta I_B$ & Amplification \\
Sature & $V_{BE} \approx 0{,}6\,\mathrm{V}$, $V_{CE} \approx 0{,}2\,\mathrm{V}$ & Limite par circuit & Interrupteur ferme \\
\bottomrule
\end{tabular}
\end{center}
En saturation, $I_C = \beta I_B$ n'est \emph{plus} valable.

\medskip\hrule\vspace{4pt}

% -- 1.3 MOSFET, IGBT et hacheur abaisseur ----------------------------
\textbf{1.3 MOSFET et IGBT - hacheur abaisseur}

\smallskip
\begin{center}
\renewcommand{\arraystretch}{1.3}
\begin{tabular}{lll}
\toprule
\textbf{Composant} & \textbf{Commande} & \textbf{Usage} \\
\midrule
MOSFET & Tension $V_{GS}$ & Haute frequence, BT \\
IGBT & Tension $V_{GE}$ (+ diode anti-parallele) & Variateurs, onduleurs HT \\
\bottomrule
\end{tabular}
\end{center}

\textbf{Hacheur abaisseur (Buck) - etats :}
\begin{itemize}[itemsep=1pt]
  \item $t \in [0\,;\,\alpha T]$ : $Q_1$ ferme, $D_1$ bloquee - $v_{Q1} = 0$, $v_{D1} = -V_{DC}$, $v_L = V_{DC} - V_s$ - $i_L$ croet
  \item $t \in [\alpha T\,;\,T]$ : $Q_1$ ouvert, $D_1$ passante (roue libre) - $v_{Q1} = V_{DC}$, $v_{D1} = 0$, $v_L = -V_s$ - $i_L$ decroet
\end{itemize}
\[
  \langle V_{sortie} \rangle = \alpha \cdot V_{DC}
  \qquad;\qquad
  \Delta i_L = \frac{V_{DC}\,\alpha\,(1-\alpha)\,T}{L}
  \qquad;\qquad
  \Delta i_{L,max} \text{ pour } \alpha = 0{,}5
\]

\end{tcolorbox}

% -----------------------------------------------------------------------
%  2. FORMULAIRE SYNTHeTIQUE
% -----------------------------------------------------------------------

\begin{tcolorbox}[
  ficheSection=colElec,
  title={2. Formulaire synthetique},
]
\small
\begin{multicols}{2}
\renewcommand{\arraystretch}{1.45}
\begin{tabular}{@{}p{4.5cm}ll@{}}
\toprule
\textbf{Formule} & \textbf{Usage} & \textbf{U.} \\
\midrule
\multicolumn{3}{@{}l}{\footnotesize\textit{Diode}} \\
$V_{seuil,Si} = 0{,}6\,\mathrm{V}$ & Seuil & V \\
$E = Ri + V_{AK}$ & Droite de charge & - \\
\midrule
\multicolumn{3}{@{}l}{\footnotesize\textit{BJT}} \\
$I_C = \beta I_B$ (regime actif) & Amplification & A \\
$V_{CE,sat} \approx 0{,}2\,\mathrm{V}$ & Saturation & V \\
\midrule
\multicolumn{3}{@{}l}{\footnotesize\textit{Hacheur}} \\
$\langle V_s \rangle = \alpha V_{DC}$ & Tension moy. & V \\
$\Delta i_L = V_{DC}\alpha(1-\alpha)T/L$ & Ondulation & A \\
\bottomrule
\end{tabular}

\columnbreak

\renewcommand{\arraystretch}{1.45}
\begin{tabular}{@{}p{4.5cm}ll@{}}
\toprule
\textbf{Formule} & \textbf{Usage} & \textbf{U.} \\
\midrule
\multicolumn{3}{@{}l}{\footnotesize\textit{Courant inductance}} \\
$i_L(t) = i_0 + (V_{DC}-V_s)t/L$ & $Q_1$ ferme & A \\
$i_L(t) = i_{max} - V_s(t-\alpha T)/L$ & $D_1$ passante & A \\
\midrule
\multicolumn{3}{@{}l}{\footnotesize\textit{Inductance min.}} \\
$L_{min} = V_{DC}\alpha(1-\alpha)T/\Delta i_{max}$ & Contin. condu. & H \\
\midrule
\multicolumn{3}{@{}l}{\footnotesize\textit{Donnees typiques}} \\
$V_{DC} = 750\,\mathrm{V}$ (traction 2019) & Bus DC & V \\
$V_{DC} = 48\,\mathrm{V}$ (Bakus 2023) & Bus DC & V \\
$\beta \approx 50$$500$ & BJT & - \\
\bottomrule
\end{tabular}
\end{multicols}
\normalsize
\end{tcolorbox}

% -----------------------------------------------------------------------
%  3. QUESTIONS TYPES
% -----------------------------------------------------------------------

\begin{tcolorbox}[
  ficheSection=colElec,
  title={3. Questions types - extraites des sujets reels},
]

% -- QT1 -------------------------------------------------------------

% Mini-procedure de resolution type copie concours
\begin{tcolorbox}[questiontype,title={Procedure type copie concours (5 etapes)}]
\begin{enumerate}[leftmargin=1.5em]
  \item \textbf{Qualifier le modele} : hypotheses, domaine de validite, grandeurs d'entree/sortie.
  \item \textbf{Ecrire les lois} : equations de base avec symboles et unites SI.
  \item \textbf{Isoler l'inconnue} : expression litterale, puis application numerique tracee.
  \item \textbf{Controler l'ordre de grandeur} : coherence physique et unite finale.
  \item \textbf{Conclure concours} : phrase courte orientee decision technique.
\end{enumerate}
\end{tcolorbox}

% Visuel de methode : chemin de resolution CAPET
\begin{center}
\begin{tikzpicture}[font=\small, node distance=7mm and 6mm, >=Stealth]
  \node[draw, rounded corners, fill=colElec!8, inner sep=4pt] (n1) {Notions};
  \node[draw, rounded corners, fill=colElec!8, right=of n1, inner sep=4pt] (n2) {Methode};
  \node[draw, rounded corners, fill=colElec!8, right=of n2, inner sep=4pt] (n3) {Application annale};
  \node[draw, rounded corners, fill=colElec!8, right=of n3, inner sep=4pt] (n4) {Verification jury};
  \draw[->, thick, colElec] (n1) -- (n2);
  \draw[->, thick, colElec] (n2) -- (n3);
  \draw[->, thick, colElec] (n3) -- (n4);
\end{tikzpicture}
\end{center}
\begin{tcolorbox}[questiontype,
  title={QT1 - Courant dans l'inductance d'un hacheur abaisseur \hfill\normalfont\footnotesize E2 disc. 2025 (Q27-Q28)}
]
\textbf{enonce.}
Hacheur abaisseur avec $Q_1$ (IGBT) et $D_1$ (roue libre). $V_{DC} = 48\,\mathrm{V}$, $V_s = 24\,\mathrm{V}$, $L = 2\,\mathrm{mH}$, $\alpha = 0{,}6$, $T = 100\,\mu\mathrm{s}$.
Exprimer $i_{L1}(t)$ sur chaque intervalle par les lois de Kirchhoff.

\tcbline

\textbf{Corrige.}

\textbf{Intervalle 1 - $t \in [0\,;\,\alpha T]$ ($Q_1$ ferme, $D_1$ bloquee) :}\\
Maille : $V_{DC} = v_L + V_s \Rightarrow v_L = V_{DC} - V_s = 48 - 24 = 24\,\mathrm{V}$
\[
  v_L = L\,\frac{di_L}{dt} \;\Rightarrow\; i_L(t) = i_0 + \frac{V_{DC} - V_s}{L}\,t = i_0 + \frac{24}{2\times10^{-3}}\,t \quad\text{(croissant)}
\]

\textbf{Intervalle 2 - $t \in [\alpha T\,;\,T]$ ($Q_1$ ouvert, $D_1$ passante) :}\\
Maille : $0 = v_L + V_s \Rightarrow v_L = -V_s = -24\,\mathrm{V}$
\[
  i_L(t) = i_{max} - \frac{V_s}{L}\,(t - \alpha T) = i_{max} - \frac{24}{2\times10^{-3}}\,(t - 60\,\mu\mathrm{s}) \quad\text{(decroissant)}
\]
Ondulation : $\Delta i_L = \frac{48 \times 0{,}6 \times (1-0{,}6) \times 100\times10^{-6}}{2\times10^{-3}} = \mathbf{0{,}576\,\mathrm{A}}$
\end{tcolorbox}

\vspace{4pt}

% -- QT2 -------------------------------------------------------------

\begin{tcolorbox}[questiontype,
  title={QT2 - Rapport cyclique d'ondulation maximale et inductance minimale \hfill\normalfont\footnotesize E2 disc. 2023 (Q C6, robot Bakus)}
]
\textbf{enonce.}
Hacheur abaisseur : $V_{DC} = 48\,\mathrm{V}$, $T = 100\,\mu\mathrm{s}$, ondulation maximale admissible $\Delta i_{max} = 4\,\mathrm{A}$.
(a) Trouver le rapport cyclique $\alpha$ qui maximise $\Delta i_L$. (b) Calculer l'inductance minimale $L_{min}$.

\tcbline

\textbf{Corrige.}

\textbf{(a)} $\Delta i_L = V_{DC}\,\alpha(1-\alpha)T/L$. Maximiser $\alpha(1-\alpha)$ - deriver par rapport - $\alpha$ et annuler :
\[
  \frac{d}{d\alpha}[\alpha(1-\alpha)] = 1 - 2\alpha = 0 \;\Rightarrow\; \alpha^* = \mathbf{0{,}5}
\]
\textbf{(b)} En $\alpha = 0{,}5$ :
\[
  L_{min} = \frac{V_{DC} \times 0{,}5 \times 0{,}5 \times T}{\Delta i_{max}} = \frac{48 \times 0{,}25 \times 100 \times 10^{-6}}{4} = \frac{1{,}2 \times 10^{-3}}{4} = \mathbf{0{,}3\,\mathrm{mH}}
\]
\textit{En pratique, on choisit $L > L_{min}$ pour toutes les valeurs de $\alpha$.}
\end{tcolorbox}

\vspace{4pt}

% -- QT3 -------------------------------------------------------------

\begin{tcolorbox}[questiontype,
  title={QT3 - Diode de protection - entree analogique d'un CAN \hfill\normalfont\footnotesize E2 appl. 2022 (Q27)}
]
\textbf{enonce.}
Une diode $D_1$ (Si) est connectee entre l'entree analogique du CAN et la tension de reference $V_{ref}$. Expliquer son rele de protection.

\tcbline

\textbf{Corrige.}

La diode $D_1$ est montee anode cete signal, cathode cete $V_{ref}$.

En fonctionnement normal ($V_{entree} < V_{ref}$) : $V_{AK} = V_{entree} - V_{ref} < 0$ - diode \textbf{bloquee}, aucune influence sur le signal.

En cas de surtension ($V_{entree} > V_{ref} + V_{seuil}$) : $V_{AK} > 0{,}6\,\mathrm{V}$ - diode \textbf{passante} - elle ecrete la tension - :
\[
  V_{entree,max} = V_{ref} + V_{seuil} = V_{ref} + 0{,}6\,\mathrm{V}
\]
La surtension est absorbee par la diode et ne penetre pas dans l'entree du CAN, evitant sa destruction.
\end{tcolorbox}

\end{tcolorbox}

% -----------------------------------------------------------------------
%  4. ATTENTES DU JURY
% -----------------------------------------------------------------------

\begin{tcolorbox}[
  jury,
  title={4. Attentes du jury - Erreurs recurrentes et points de vigilance},
]

\begin{enumerate}[label=\textbf{\arabic*.}, leftmargin=*, itemsep=5pt]

  \item \textbf{Oublier la diode de roue libre.}
  Sans $D_1$ en anti-parallele avec la charge inductive, la tension aux bornes du transistor diverge - l'ouverture ($v = L\,di/dt \to \infty$) et le detruit. Toujours justifier sa presence.

  \item \textbf{Ondulation maximale - $\alpha = 0{,}5$.}
  $\Delta i_L \propto \alpha(1-\alpha)$ est maximisee pour $\alpha = 0{,}5$, pas pour $\alpha = 1$. Erreur classique dans la formule.

  \item \textbf{Rapport cyclique et tension moyenne.}
  $\langle V_s \rangle = \alpha V_{DC}$ n'est valable qu'en \emph{conduction continue} ($i_L > 0$ en permanence). En conduction discontinue, la relation change.

  \item \textbf{Saturation du BJT $\neq$ regime actif.}
  En saturation, $I_C = \beta I_B$ n'est plus valable et $V_{CE} \approx 0{,}2\,\mathrm{V}$. Le transistor se comporte comme un interrupteur ferme.

  \item \textbf{Signe de $V_{AK}$ pour la diode.}
  En passant, $V_{AK} > 0$. En bloquee, $V_{AK} < 0$. L'erreur de signe est la faute la plus frequente.

  \item \textbf{Schemas equivalents pour chaque intervalle.}
  Le jury attend que le candidat dessine le circuit equivalent ($Q_1$ ouvert ou ferme, $D_1$ passante ou bloquee) avant d'appliquer les lois de Kirchhoff.

\end{enumerate}

\medskip
\begin{tcolorbox}[
  colback=colJuryFrame!12!colJury,
  colframe=colJuryFrame,
  left=4pt, right=4pt, top=2pt, bottom=2pt,
  fontupper=\small\itshape
]
Rapport jury 2024 : "Le jury deplore le manque de maetrise des lois fondamentales (loi des neuds, loi des mailles) pour les composants elementaires."\\
Rapport jury 2025 : "Les resultats donnes directement, sans demarche et sans justification ne sont pas valorises. Dessiner les schemas equivalents pour chaque intervalle de commutation."
\end{tcolorbox}

\end{tcolorbox}

% -----------------------------------------------------------------------
%  PIED - PReSENCE DANS LES SUJETS
% -----------------------------------------------------------------------

\begin{tcolorbox}[
  enhanced, breakable,
  colback=colGrayLight,
  colframe=colSeparator,
  left=4pt, right=4pt, top=3pt, bottom=3pt,
  before skip=4pt, after skip=0pt,
  fonttitle=\small\bfseries,
  title={Presence dans les sujets analyses (2013-2026)},
]
\begin{center}
\footnotesize\renewcommand{\arraystretch}{1.15}
\begin{tabular}{lccccccccccccc}
\toprule
\textbf{Annee} & 13 & 14 & 15 & 16 & 17 & 18 & 19 & 20 & 21 & 22 & 23 & 24 & 25 \\
\midrule
E1 (commune)   & $\checkmark$ & - & - & - & - & - & - & - & - & - & - & - & $\checkmark$ \\
E2 discipl. IE & - & - & - & - & - & - & - & - & - & $\checkmark$ & $\checkmark$ & - & $\checkmark$ \\
\bottomrule
\end{tabular}
\end{center}
\smallskip
\footnotesize\textbf{Sous-themes frequents :}
hacheur abaisseur - courant et ondulation dans $L$ (3/13) - rapport cyclique et tension moyenne (4/13) - diode de roue libre (4/13) - IGBT/MOSFET en convertisseurs (3/13) - diode de protection CAN (2/13)

\medskip
\textbf{Ordres de grandeur (reperes rapides).}
\begin{itemize}[leftmargin=1.3em]
  \item Reseau BT FR : tension phase-neutre typique \SI{230}{\volt}, tension phase-phase typique \SI{400}{\volt}, frequence \SI{50}{\hertz}.
  \item Courants d'etude frequents en sujet : de quelques \si{\ampere} a quelques dizaines de \si{\ampere}.
  \item En conversion d'energie, un rendement attendu reste generalement entre 0{,}8 et 0{,}98 selon la technologie.
\end{itemize}

\medskip
\textbf{Synthese express (a refaire sans support).}
\begin{enumerate}[leftmargin=1.4em]
  \item Poser les hypotheses du modele et les conventions de signe.
  \item Ecrire 2 a 4 relations cles avec unites.
  \item Realiser un calcul numerique complet et verifier l'ordre de grandeur.
  \item Justifier le resultat par un argument technique de niveau CAPET.
\end{enumerate}

\end{tcolorbox}

\CAPETENDFILE









